\subsection{The fixed-effect comparison}
Table~\ref{tab:fe} compares four specifications on the same 486,655 complete-case observations. All include contemporaneous volume, capitalization, and the daily price range. The security/month coefficient is 0.0041. Adding daily effects yields 0.0145; market-day effects yield 0.0129. The sign changes to -0.0033 when sector-month effects are included. This sign change occurs without changing the observations or financial tail caps. It is therefore a specification difference, rather than a consequence of comparing different complete-case samples. Figure~\ref{fig:specifications} displays the estimates and uncertainty.

The positive local-market-day association and its disappearance after sector-month controls suggest that slower sector-level co-movement matters to the monthly result. This comparison does not identify whether that component reflects confounding sector conditions or a persistent information channel absorbed by the sector-month effects. It establishes sensitivity of the descriptive coefficient, rather than proving that the more saturated model isolates a causal effect.

\subsection{Principal estimate and economic scale}
Table~\ref{tab:primary} reports the model with prior-session controls. The coefficient is -0.00143, with sector/date standard error 0.00220 and an approximate 95\% interval [-0.00650, 0.00364]. It is not distinguishable from zero at conventional levels ($p=0.533$). Scaling by the sample standard deviation of content gives -0.102 basis points, with an interval [-0.462, 0.258], compared with a winsorized mean spread of 47.29 basis points. These are conditional associations in the fitted sample, not the estimated consequences of an intervention.

Security/date clustering gives $p=0.188$, while date-score HAC gives $p=0.521$ at bandwidth 20. The sector score-sign diagnostic gives 0.512. These calculations use distinct assumptions and cannot repair endogeneity, but none supplies evidence for a positive coefficient in this model. The supplement reports all three HAC bandwidths and the exposure-variation contribution of each sector.

\subsection{Timing and sensitivity}
Table~\ref{tab:timing} compares current content, previous-session content, previous-calendar-day content, and a future-session lead on a common sample. The previous-session estimate on the larger available timing sample is -0.0010 ($p=0.652$). The future-session lead in the common sample is 0.0078 ($p=0.135$). The lead comparison does not establish reverse causality or validate the lag; persistent information conditions and serially correlated content can generate both. The estimates do not establish a distinctive positive delayed response after the local market and sector controls.

The supplementary tables report transformations of the spread, alternative content measures, the matched-post-only sample, joint adjustment for post volume and relevance, and sector/market leave-out checks. Raw-spread and log-spread coefficients do not have a common numerical scale, and their signs must not be selected to support a preferred mechanism. The formal post-2019 slope interaction is -0.0019 ($p=0.687$); it is a direct test of a slope difference, rather than a comparison of whether two separate subperiod probabilities cross 0.05. All these checks are descriptive sensitivity analyses, with no independent causal interpretation.
